\section{T-I Debugging and T-II Status}
\label{sec:debugging}

\subsection{Why the first Push-T run was rejected}

An earlier spatial RGB checkpoint reached 0\% closed-loop success despite low offline diffusion
loss, high cosine similarity to stored demonstration actions, and strong sensitivity to shuffled
RGB. Those diagnostics showed that the model fit its dataset; they did not establish that the
dataset represented the intended observation-action relation. Treating that checkpoint as a valid
base policy would have made every downstream failure frequency scientifically ambiguous.

The replay audit identified a concrete temporal defect. In a state-restored trajectory, action
$a_t$ advances the simulator from $s_t$ to $s_{t+1}$. The subsequent rendered observation must
therefore correspond to $s_{t+1}$. The original GPU replay restored $s_t$ after the action,
creating a shifted supervised mapping while preserving plausible images and valid tensor shapes.
Section~\ref{sec:aligned-replay} describes the fail-closed correction and validation gates. The
pre-fix checkpoint was not resumed because changing the dataset changes the learned conditional
mapping.

\noindent\textbf{Fail-closed alignment patch (abridged).}
\begin{minipage}{\linewidth}
\scriptsize
\begin{verbatim}
OLD = "set_state_dict(env_states_batch[t])"
NEW = "set_state_dict(env_states_batch[t + 1])"
if source.count(OLD) != 1:
    raise RuntimeError("inspect installed source")
source = source.replace(OLD, NEW)
\end{verbatim}
\end{minipage}

\subsection{Controls and independent troubleshooting}

Several controls separated data, architecture, and infrastructure failures. First, a spatial
feature encoder replaced global visual pooling because Push-T withholds privileged object and goal
pose; absolute image location is therefore task-relevant. Second, an RGB PushCube control reached
82.0\% $\pm$ 4.6 percentage points over 300 rollouts, validating the common image replay, policy,
EMA, and evaluator. Third, the corrected Push-T replay matched the source controller and original
1,024-environment collection parallelism rather than changing several variables simultaneously.

Infrastructure problems were also made explicit rather than hidden. Python 3.12 and headless
Vulkan were verified before training. HCE jobs call Python directly because nesting \texttt{srun}
inside the allocation caused CPU-binding failures. Runpod training was detached from SSH, logged
to persistent storage, checkpointed every 5k, and resumed from the newest validated checkpoint
after termination. These decisions are part of reproducibility, not incidental operational detail.

\subsection{Limitations and hidden assumptions}

The result supports a trained visual Push-T baseline, but it does not establish robustness:
\begin{itemize}
    \item Only one policy-training seed (seed 1) was affordable. The reported standard deviation
    measures environment-seed variation for one fixed policy, not training-seed variability.
    \item Checkpoint selection repeatedly evaluates 20 episodes. Its maximum is statistically
    noisy and selection-biased, as demonstrated by the lower 300-rollout estimate.
    \item Push-T GPU replay is sensitive to collection parallelism and simulator state restoration.
    Matching 1,024 environments reduces a known discrepancy but does not prove bitwise equivalence
    across hardware or ManiSkill versions.
    \item \successonce is the official baseline metric, but it credits transient goal entry. We
    retain \successend, maximum overlap, final overlap, and videos to expose this distinction.
    \item The loss trace starts at 10.1k because the original Colab TensorBoard event file was not
    preserved. Checkpoints and evaluation metrics from that segment remain available.
\end{itemize}

These caveats motivate T-II's pose-conditioned evaluation: aggregate success alone cannot reveal
which initial configurations fail, whether distinct mechanisms are reproducible, or whether a
later residual policy improves targeted failures without harming nominal behavior.

\subsection{T-II: implemented protocol, incomplete experiments}

T-II was not completed. We implemented a separate \texttt{t-ii/} package that samples exact
subsets of Push-T's nominal initial-pose domain, records the initial T pose and TCP pose, saves
per-step overlap trajectories, assigns non-exclusive screening tags, and aggregates success with
Wilson intervals. Unit tests cover wrapped angular ranges, exact pose-conditioned resets, metric
serialization, and failure-tag logic. These tests establish the evaluation machinery, not two
empirical failure modes.

The preregistered discovery protocol was 200 full-distribution rollouts from relative block
position $x\in[-0.10,0.10]$ m, $y\in[-0.10,0.20]$ m and yaw $\theta\in[0,360)^\circ$. Candidate
regions would be selected only after jointly inspecting initial-pose bins, overlap trajectories,
and videos. Two useful screening hypotheses were (i) low progress after loss of contact and
(ii) progress regression after transient partial or successful goal overlap. These are plausible
mechanisms, not reported findings. Each accepted mode would then require a representative MP4 and
100 targeted rollouts at each seed in $\{0,1,2\}$.

The valid visual run could not be executed before the deadline. The original A40 completed T-I
and was safely backed up. A replacement A100 booted and received the hash-verified checkpoint, but
RGB environment creation failed with \texttt{vk::createInstanceUnique:
ErrorIncompatibleDriver}; graphics-capable A40, RTX A6000/4090, RTX 4000 Ada, and L4 instances
were unavailable. A local Apple M4 smoke test reproduced the missing-Vulkan barrier. We rejected
a state-only substitute because it would not evaluate the trained RGB policy. Thus there are no
defensible T-II per-mode rates or videos, and none are fabricated or inferred from T-I diagnostic
rollouts.

We also avoid treating an external number as a required target. ManiSkill issue \#882 reports one
state-input reproduction near 0.15--0.20 and asks why a paper reports 0.4--0.5, but it does not
establish an expected RGB result or resolve differences in data and simulation
settings~\cite{maniskill_issue882}. Our baseline is therefore defined by its disclosed protocol,
not by post-hoc agreement with that thread.
